\hypertarget{defter}{%
\chapter{Ek A --- Kümülatif defter}\label{defter}}

Her NACE grubu tamamlandığında güncellenir. Toplamlar Faz 6'daki makro mutabakatta kullanılır.

\textbf{Son güncelleme:} Faz 3 parça 2 (13 Tekstil, 14 Giyim, 15 Deri), 18 Ekim 2026. Değerler 10--15'in kümülatif toplamı (54 sınıf). Sınıf düzeyi değerler \kod{veri/\allowbreak{}nace\_\allowbreak{}matris.\allowbreak{}csv}'den; parça ayrıntısı \belgebag{nace_kartlar_13-15}{}{\kod{nace\_\allowbreak{}kartlar\_\allowbreak{}13-\allowbreak{}15.\allowbreak{}md}} ve \belgebag{nace_kartlar_10-12}{}{\kod{nace\_\allowbreak{}kartlar\_\allowbreak{}10-\allowbreak{}12.\allowbreak{}md}} Bölüm 6. Grup toplamlarının yüzdelikleri çekilişlerin toplamından (V-38).

\begingroup\gbtablo\small

\begin{longtable}[]{@{}
  >{\raggedright\arraybackslash}p{(\columnwidth - 4\tabcolsep) * \real{0.1592}}
  >{\raggedright\arraybackslash}p{(\columnwidth - 4\tabcolsep) * \real{0.6566}}
  >{\raggedright\arraybackslash}p{(\columnwidth - 4\tabcolsep) * \real{0.1841}}@{}}
\toprule\noalign{}
\begin{minipage}[b]{\linewidth}\raggedright
\textbf{Kalem}
\end{minipage} & \begin{minipage}[b]{\linewidth}\raggedright
\textbf{Değer}
\end{minipage} & \begin{minipage}[b]{\linewidth}\raggedright
\textbf{Not}
\end{minipage} \\
\midrule\noalign{}
\endhead
\bottomrule\noalign{}
\endlastfoot
Tamamlanan / kalan sınıf sayısı & 54 / 176 C sınıfı (10--12 ve 13--15 tamam); 0 / 29 ek kapsam satırı (B 15, D 1, E 2, A 3, hizmet 8) & Kapsam listesi: \kod{veri/\allowbreak{}nace\_\allowbreak{}kapsam.\allowbreak{}csv} \\
Kullanılan / kalan enflasyon bütçesi (bp) & Birinci tur, 10--15. Çıkış şokları (önlem değil): S2 +2,6, S4 +21,8 (c2; 13--15'in payı 15,6). Önlemler: S2 yolu −0,11, S4 yolu −13,56 (13--15'te giyim ve ayakkabıda KDV takası −12,0). Net: S2 2,4 / 50; S4 8,2 / 50 (yıllık tavan 30). İkinci tur dahil Faz 6'da & V-05, V-43; \kod{nace\_\allowbreak{}kartlar\_\allowbreak{}13-\allowbreak{}15.\allowbreak{}md} Tablo 6.1 \\
Kümülatif cari denge katkısı (milyar \$) & 10--15, T+5, yıllık. Senaryo P50: S2 −0,93, S4 −2,75 (S4 P10 −4,15; 1/98 askısında ortalama −3,46). Önlemler: S2 +0,22, S4 +0,01 \etiketT{}/\etiketH{} & \kod{veri/\allowbreak{}f3\_\allowbreak{}senaryo\_\allowbreak{}grup.\allowbreak{}csv}, \kod{veri/\allowbreak{}f3\_\allowbreak{}senaryo\_\allowbreak{}aski.\allowbreak{}csv} \\
Kümülatif mali maliyet (milyar TL ve GSYH \%) & 10--15: S2 yolu 1,65 mlr TL/yıl (\%0,002), S4 yolu 22,76 mlr TL/yıl (\%0,031); S4'ün 17,5'i KDV takasının net maliyeti & Tavan: GSYH'nin \%0,5'i/yıl (2026'da ≈ \$9,1 mlr; 2028'de ≈ \$10,5 mlr) \\
Riskteki istihdam & 10--15: S2 P50 −14.630; S4 P50 −40.241, P10 −62.512; 1/98 askısında ortalama −42.977. S4 P10'da kaybı \%5'i aşan sınıflar: 13.20, 13.91, 13.92, 14.13, 14.14, 14.19, 14.31, 14.39 (trikoda −\%19,2). PEM kümülasyonu kanalı S4'te −4.750. Faz 3 yöntemiyle imalatın tamamı (geçici): S2 P50 −24,3 bin (P10 −42,5); S3 −29,1 bin; S4 −46,1 bin (P10 −70,1), kaybı \%5'i aşan sınıf S4 P50'de 9, P10'da 17. Teşhis tabanı (S4, statik, telafisiz): 207 bin brüt & \kod{veri/\allowbreak{}f3\_\allowbreak{}senaryo\_\allowbreak{}ozet.\allowbreak{}csv}; \kod{01\_\allowbreak{}teshis.\allowbreak{}md} §4.1 \\
Açık kalan veri boşlukları & VB-02, 04, 06, 08--12, 14--17, 21--27, 30--39 (kapanan: VB-01, 03, 05, 07, 13, 18). VB-32 13--15'te bağlayıcı değil & \kod{01\_\allowbreak{}teshis.\allowbreak{}md} §8; \kod{nace\_\allowbreak{}kartlar\_\allowbreak{}13-\allowbreak{}15.\allowbreak{}md} Ek E \\
Varsayım kaydına eklenen yeni maddeler & V-01 \ldots{} V-43 (Faz 3: V-30--V-39 senaryo kalibrasyonu, 1/98 askı olasılığı, ProTEE σ, sınır sürtünmesi, arketip eşikleri, SS, CKE, KN8 eşlemesi, grup yüzdelikleri, eylem takvimi; parça 2: V-40 PEM kümülasyonu kanalı, V-41 tekstil ve giyim menşe parametreleri, V-42 geri ödeme yasağı, V-43 KDV takası) & \kod{00\_\allowbreak{}plan.\allowbreak{}md} §3 \\
\end{longtable}

\endgroup

\hypertarget{defter-teux15fhis-tabanux131-faz-1}{%
\section{Teşhis tabanı (Faz 1)}\label{defter-teux15fhis-tabanux131-faz-1}}

Sınıf kartlarının ve makro köprünün karşılaştırılacağı başlangıç değerleri.

\begingroup\gbtablo\small

\begin{longtable}[]{@{}
  >{\raggedright\arraybackslash}p{(\columnwidth - 4\tabcolsep) * \real{0.3385}}
  >{\raggedright\arraybackslash}p{(\columnwidth - 4\tabcolsep) * \real{0.4485}}
  >{\raggedright\arraybackslash}p{(\columnwidth - 4\tabcolsep) * \real{0.2130}}@{}}
\toprule\noalign{}
\begin{minipage}[b]{\linewidth}\raggedright
\textbf{Gösterge}
\end{minipage} & \begin{minipage}[b]{\linewidth}\raggedright
\textbf{Değer}
\end{minipage} & \begin{minipage}[b]{\linewidth}\raggedright
\textbf{Kaynak ve etiket}
\end{minipage} \\
\midrule\noalign{}
\endhead
\bottomrule\noalign{}
\endlastfoot
Cari denge 2025 & −\$30,3 mlr (−\%1,9 GSYH) & \etiketV{} EVDS3 \\
Altın ve enerji hariç (çekirdek) cari denge 2025 & +\$37,1 mlr (+\%2,3) & \etiketV{} EVDS3 \\
Cari denge, 12 ay (Eyl 2025--Ağu 2026) & −\$42,1 mlr; çekirdek +\$26,0 mlr & \etiketV{} EVDS3 \\
OVP cari denge patikası 2027 / 2028 / 2029 & −\%1,9 / −\%1,8 / −\%1,6 & \etiketV{} OVP Tablo 1.1 \\
T+5 cari denge hedef boşluğu (2033 / 2034) & \$28--45 mlr / \$25--48 mlr & \etiketH{} \kod{01\_\allowbreak{}teshis.\allowbreak{}md} §3.6 \\
TÜFE yıl sonu 2026 / 2027 / 2028 / 2029 & \%28,4 / \%21,0 / \%13,5 / \%9,0 & \etiketV{} OVP \\
İmalat istihdamı & YSHİ 5,03 mn (2024); SGK 4/a 4,03 mn (2025) & \etiketV{} \\
AB ihracatına doğrudan bağlı imalat istihdamı & 858 bin (YSHİ bazı); 675 bin (SGK bazı) & \etiketT{} §4 \\
S4 statik ihracat kaybı (σ = 3 / 5 / 8) & −€13,8 / −21,2 / −30,4 mlr & \etiketT{} Faz 0 S4 hesabı \\
S4 TÜFE birinci tur (pasif): yalnızca MFN / MFN + İGV & 32--42--53 bp / 59--78--98 bp (girdi-çıktı modeli; ilk kaba hesap 21--40 / 46--86 bp) & \etiketT{} \kod{veri/\allowbreak{}politika\_\allowbreak{}io\_\allowbreak{}s4\_\allowbreak{}tufe.\allowbreak{}csv} \\
STA asimetrisi: Türk ihracatının ödediği MFN & \$1,2 mlr (yürürlükte) + \$0,3 mlr (bekleyen) & \etiketT{} §2.1 \\
\end{longtable}

\endgroup

\hypertarget{defter-politika-uzayux131-tabanux131-faz-2a}{%
\section{Politika uzayı tabanı (Faz 2A)}\label{defter-politika-uzayux131-tabanux131-faz-2a}}

Faz 3--6'daki kararların karşılaştırılacağı referans değerler. Ayrıntı: \belgebag{02a_politika_uzayi}{}{\kod{02a\_\allowbreak{}politika\_\allowbreak{}uzayi.\allowbreak{}md}}.

\begingroup\gbtablo\small

\begin{longtable}[]{@{}
  >{\raggedright\arraybackslash}p{(\columnwidth - 4\tabcolsep) * \real{0.2435}}
  >{\raggedright\arraybackslash}p{(\columnwidth - 4\tabcolsep) * \real{0.5318}}
  >{\raggedright\arraybackslash}p{(\columnwidth - 4\tabcolsep) * \real{0.2246}}@{}}
\toprule\noalign{}
\begin{minipage}[b]{\linewidth}\raggedright
\textbf{Gösterge}
\end{minipage} & \begin{minipage}[b]{\linewidth}\raggedright
\textbf{Değer}
\end{minipage} & \begin{minipage}[b]{\linewidth}\raggedright
\textbf{Kaynak ve etiket}
\end{minipage} \\
\midrule\noalign{}
\endhead
\bottomrule\noalign{}
\endlastfoot
S4 TÜFE birinci tur, pasif (MFN + İGV) & 59 / 78 / 98 bp (ρ = 0,6 / 0,8 / 1,0); κ = 0,5 ile ≈182 bp & \etiketT{} \kod{veri/\allowbreak{}politika\_\allowbreak{}s4\_\allowbreak{}aktif.\allowbreak{}csv} \\
Aktif S4 önerilen paket: c2 + P1 + geniş takas & TÜFE 2,5 / 3,3 / 4,1 bp; net mali +\%0,05 GSYH (tahsil edilebilir); takas maliyeti \$5,7 mlr; koruma kaybı 0 & \etiketT{} §4.1 \\
Aktif S4 alternatifleri (ρ = 0,8) & c + dar 7,5 bp, −\%0,02; a + geniş 4,0 bp, +\%0,07; b 1,5 bp, −\%0,04 ve \$2,9 mlr koruma kaybı & \etiketT{} §4.1 \\
Vergi takası, on grup (T-A, rejim c2, merkez) & TÜFE 0 bp (pasif +47); net mali −\$0,60 mlr (−0,71 ile +0,13); istihdam +31 bin (18--36). Rejim c'de pasif +35 bp, takas −\$0,46 mlr & \etiketT{} \kod{veri/\allowbreak{}politika\_\allowbreak{}takas\_\allowbreak{}ozet.\allowbreak{}csv} \\
Otomobil takası & ÖTV \%65 → \%56,2; AB +\%4,1, yerli −\%5,4; net mali −\$0,15 mlr; T-0 (yalnızca yerli) eşdeğeri −\$3,0 mlr ve DTÖ'ye aykırı & \etiketT{} §1.5 \\
S4 net imalat istihdamı (σ = 5) & Yalnızca MFN −117 bin; c2 −48 bin; c2 + takas ≈ −35 bin; giyim −58,7 bin & \etiketT{} \kod{veri/\allowbreak{}politika\_\allowbreak{}s4\_\allowbreak{}net\_\allowbreak{}istihdam\_\allowbreak{}ozet.\allowbreak{}csv} \\
Gümrük vergisi tahsilat oranı & Toplam \%50 (2025: \$8,15 mlr gerçekleşme / \$16,2 mlr yasal); tüketim ve karma mallarda \%100, diğerlerinde \%38 & \etiketV{}/\etiketT{} §3.4 \\
AB'nin S4 ihracat kaybı (σ = 5) & b \$16,4 mlr; c \$18,7 mlr; c2 \$28,5 mlr; a \$31,1 mlr & \etiketT{} §4.1 \\
İhracat tarafı tam telafi maliyeti & Yılda ≈\$38 mlr (2028 GSYH'sinin \%1,8'i); telafi edici vergiye açık & \etiketT{} §5 \\
Zaman kritik DTÖ eylemi & XXVIII:5 çekincesi, son tarih 31.12.2026; DS629 DSB, 27.10.2026 & \etiketV{} §7 \\
\end{longtable}

\endgroup

\hypertarget{defter-seuxe7enekler-tabanux131-faz-2b}{%
\section{Seçenekler tabanı (Faz 2B)}\label{defter-seuxe7enekler-tabanux131-faz-2b}}

S0'a (GB) göre yerleşmiş etkiler. Ayrıntı: \belgebag{02b_secenekler_hukuk}{}{\kod{02b\_\allowbreak{}secenekler\_\allowbreak{}hukuk.\allowbreak{}md}} §2.

\begingroup\gbtablo\footnotesize

\begin{longtable}[]{@{}
  >{\raggedright\arraybackslash}p{(\columnwidth - 8\tabcolsep) * \real{0.2190}}
  >{\raggedright\arraybackslash}p{(\columnwidth - 8\tabcolsep) * \real{0.1676}}
  >{\raggedleft\arraybackslash}p{(\columnwidth - 8\tabcolsep) * \real{0.2275}}
  >{\raggedright\arraybackslash}p{(\columnwidth - 8\tabcolsep) * \real{0.1680}}
  >{\raggedright\arraybackslash}p{(\columnwidth - 8\tabcolsep) * \real{0.2179}}@{}}
\toprule\noalign{}
\begin{minipage}[b]{\linewidth}\raggedright
\textbf{Gösterge}
\end{minipage} & \begin{minipage}[b]{\linewidth}\raggedright
\textbf{S2 DCFTA}
\end{minipage} & \begin{minipage}[b]{\linewidth}\raggedleft
\textbf{S3 ITA}
\end{minipage} & \begin{minipage}[b]{\linewidth}\raggedright
\textbf{S4}
\end{minipage} & \begin{minipage}[b]{\linewidth}\raggedright
\textbf{Kaynak}
\end{minipage} \\
\midrule\noalign{}
\endhead
\bottomrule\noalign{}
\endlastfoot
TR → AB ihracatı (GE, σ = 7; \%90 GA) & −\%11 (−27 ile +6) & −\%18 (ihtiyatlı) ile +\%1 (Türkiye--BK) & −\%19 (−32 ile −5) & \etiketT{} \kod{veri/\allowbreak{}politika\_\allowbreak{}cekim\_\allowbreak{}senaryolar.\allowbreak{}csv} \\
Türkiye toplam ihracatı (GE) & −\%3,3 & −\%5,3 ile +\%1,1 & −\%6,2 & \etiketT{} \\
Türkiye reel GSYH, refah vekili (GE) & −\%0,22 (−0,56 ile +0,19) & −\%0,35 ile +\%0,10 & −\%0,43 (−0,73 ile −0,10) & \etiketT{} \\
TÜFE birinci tur (bp) & +2 ile +12 (menşe maliyeti) & +2 ile +12 & 2,5--4,1 (aktif) & \etiketT{} 02a, f2e \\
Net imalat istihdamı (bin) & ≈ −33 (kaba, GE kanalı) & −54 ile +12 & ≈ −35 (aktif, sınıf düzeyi) & \etiketT{} \\
\end{longtable}

\endgroup

\hypertarget{defter-hazux131rlux131k-tabanux131-faz-2c}{%
\section{Hazırlık tabanı (Faz 2C)}\label{defter-hazux131rlux131k-tabanux131-faz-2c}}

Karar öncesi hamlelerin değer ölçüleri ve zamanlama kısıtları. Ayrıntı: \belgebag{02c_hazirlik}{}{\kod{02c\_\allowbreak{}hazirlik.\allowbreak{}md}}.

\begingroup\gbtablo\small

\begin{longtable}[]{@{}
  >{\raggedright\arraybackslash}p{(\columnwidth - 4\tabcolsep) * \real{0.2327}}
  >{\raggedright\arraybackslash}p{(\columnwidth - 4\tabcolsep) * \real{0.5745}}
  >{\raggedright\arraybackslash}p{(\columnwidth - 4\tabcolsep) * \real{0.1928}}@{}}
\toprule\noalign{}
\begin{minipage}[b]{\linewidth}\raggedright
\textbf{Gösterge}
\end{minipage} & \begin{minipage}[b]{\linewidth}\raggedright
\textbf{Değer}
\end{minipage} & \begin{minipage}[b]{\linewidth}\raggedright
\textbf{Kaynak ve etiket}
\end{minipage} \\
\midrule\noalign{}
\endhead
\bottomrule\noalign{}
\endlastfoot
Hamleler & 51: 11 A, 26 B, 14 C. Hiçbirinin çıkış olmazsa değeri negatif değil & \etiketH{} \kod{veri/\allowbreak{}hazirlik\_\allowbreak{}hamleler\_\allowbreak{}oncelik.\allowbreak{}csv} \\
Sessiz pencere (bugün → K) & T2028 ≈1 ay (K Kasım 2026); T2029 ≈5 ay (K Mart 2027) & \etiketH{} §13 \\
33 pazarla STA: S4 GSYH kaybının telafisi (σ = 5 / 7) & Pazar erişimi \%14 / \%23; karşılıklı \%42 / \%69. GB içinde GSYH +\%0,08--0,30 & \etiketT{} \kod{veri/\allowbreak{}hazirlik\_\allowbreak{}sta\_\allowbreak{}ge.\allowbreak{}csv} \\
STA'ların statik gümrük geliri kaybı & Yılda \$1,9 mlr sanayi (GSYH'nin \%0,12'si); bütün ürünler \$3,2 mlr & \etiketT{} \kod{veri/\allowbreak{}hazirlik\_\allowbreak{}sta\_\allowbreak{}tarife.\allowbreak{}csv} \\
PEM dışı ithal ara girdi payı (AB ihracatı ağırlıklı) & \%13,1 (bütünleşik üst sınır ≈\%21). Otomobil göstergesi \%36 (kural \%45) & \etiketT{} \kod{veri/\allowbreak{}hazirlik\_\allowbreak{}pem\_\allowbreak{}*.\allowbreak{}csv} \\
SKDM: AB'ye ödeme, 2028 (merkez fiyat) & Varsayılan €2.593 mn; doğrulanmış €676 mn; yurt içi fiyat AB'ye eşitse AB'ye €81 mn ve Türkiye'de €595 mn & \etiketT{} \kod{veri/\allowbreak{}hazirlik\_\allowbreak{}skdm.\allowbreak{}csv} \\
SKDM: AB'ye ödeme, 2034 & Varsayılan €3.795 mn; doğrulanmış €1.398 mn & \etiketT{} \\
SKDM aşağı akış önerisi (1.1.2028) & €10,9 mlr ihracat; yük 2028'de €245 mn (varsayılan) / €77 mn (doğrulanmış) & \etiketT{} \kod{veri/\allowbreak{}hazirlik\_\allowbreak{}skdm\_\allowbreak{}asagi.\allowbreak{}csv} \\
Net rezerv (IMF şablonu, Eylül 2026) & \$41,6 mlr; brüt \$171,6 mlr (9.10.2026) & \etiketV{}/\etiketT{} \kod{veri/\allowbreak{}hazirlik\_\allowbreak{}makro\_\allowbreak{}rezerv.\allowbreak{}csv} \\
İlan şokunun 12 aylık rezerv etkisi & Kurallı \$30--48 mlr; düzensiz \$67--90 mlr & \etiketH{} \kod{veri/\allowbreak{}hazirlik\_\allowbreak{}makro\_\allowbreak{}tampon.\allowbreak{}csv} \\
Gerekli net rezerv (makro kapı) & Kurallı ilan \$68 mlr; düzensiz \$110 mlr. Açık \$27 / \$69 mlr & \etiketT{} \kod{veri/\allowbreak{}hazirlik\_\allowbreak{}makro\_\allowbreak{}kapi.\allowbreak{}csv} \\
Zaman kritik tarihler & XXVIII:5 31.12.2026; DS629 kabul ya da temyiz 24.11.2026 (WT/DS629/6; 02a'daki 27.10.2026 uzatıldı); ETS izleme planları 27.10.2026; ilk SKDM beyanı 30.9.2027 & \etiketV{} §1.3 \\
\end{longtable}

\endgroup

\hypertarget{defter-faz-3-tabanux131-paruxe7a-1-1012}{%
\section{Faz 3 tabanı (parça 1: 10--12)}\label{defter-faz-3-tabanux131-paruxe7a-1-1012}}

Sınıf kartlarının toplamları. Faz 3 yöntemi: toplamlar Faz 2B genel dengesine kalibre, dağılım kısmi denge (\kod{kod/\allowbreak{}f3e\_\allowbreak{}senaryo.\allowbreak{}py}). Ayrıntı: \belgebag{nace_kartlar_10-12}{}{\kod{nace\_\allowbreak{}kartlar\_\allowbreak{}10-\allowbreak{}12.\allowbreak{}md}}.

\begingroup\gbtablo\small

\begin{longtable}[]{@{}
  >{\raggedright\arraybackslash}p{(\columnwidth - 4\tabcolsep) * \real{0.2350}}
  >{\raggedright\arraybackslash}p{(\columnwidth - 4\tabcolsep) * \real{0.5453}}
  >{\raggedright\arraybackslash}p{(\columnwidth - 4\tabcolsep) * \real{0.2198}}@{}}
\toprule\noalign{}
\begin{minipage}[b]{\linewidth}\raggedright
\textbf{Gösterge}
\end{minipage} & \begin{minipage}[b]{\linewidth}\raggedright
\textbf{Değer}
\end{minipage} & \begin{minipage}[b]{\linewidth}\raggedright
\textbf{Kaynak ve etiket}
\end{minipage} \\
\midrule\noalign{}
\endhead
\bottomrule\noalign{}
\endlastfoot
10--12 ağırlığı & İstihdam 636 bin (imalatın \%12,6'sı); KD \$23,3 mlr; ihracat \$24,1 mlr (\%9,6); AB'ye ihracat \$6,2 mlr (\%5,9) & \etiketV{} YSHİ 2024, BACI 2024 \\
AB'ye ihracatın rejimi (10--12) & Tarım (1/98) \%75, İTÜ \%24, GB \%1 & \etiketV{}/\etiketT{} Comext 2025 \\
S4'te AB tarifesi artışı (rota A / 1/98 kaybıyla) & \%3,3 / \%11,5 (10--12); imalat \%5,6 / \%6,0 & \etiketT{} \kod{veri/\allowbreak{}f3\_\allowbreak{}ab\_\allowbreak{}tarife\_\allowbreak{}sinif.\allowbreak{}csv} \\
10--12, bölüm toplamı P50 (P10 / P90), T+5 & S2: AB'ye ihracat −\$202 mn (−446 / +125), istihdam −299, cari denge −\$44 mn. S4: AB'ye ihracat −\$406 mn (−1.840 / −171), istihdam −106 (−3.982 / +258), cari denge +\$8 mn (−934 / +104) & \etiketT{} \kod{veri/\allowbreak{}f3\_\allowbreak{}senaryo\_\allowbreak{}grup.\allowbreak{}csv} (V-38) \\
10--12, S4'te 1/98 askısı (koşullu ortalama; olasılık 0,25) & AB'ye ihracat −\$1,72 mlr (askı olmadan −\$0,34 mlr), cari denge −\$0,86 mlr, istihdam −3,7 bin (P90 / P10: −2,0 / −5,7 bin); askıdan doğan ek kaybın \%77'si 10.39'da & \etiketT{} \kod{veri/\allowbreak{}f3\_\allowbreak{}senaryo\_\allowbreak{}aski.\allowbreak{}csv} \\
İmalat (230 sınıf), AB'ye ihracat değişimi P50 (P10 / P90) & S2 −\%11,0 (−23,4 / +2,0); S3 −\%14,7 (−26,2 / −2,6); S4 −\%19,1 (−29,7 / −7,9) & \etiketT{} GE kalibrasyonu \\
İmalat, cari denge P50 (P10 / P90), T+5, yıllık & S2 −\$2,7 mlr (−4,2 / −2,0); S3 −\$3,4 mlr; S4 −\$4,7 mlr (−6,9 / −3,4); S4'e PEM kümülasyonu kanalı dahil & \etiketT{} geçici \\
İmalat, birinci tur TÜFE (c2 ve sürtünme), P50 & S2 10,9 bp; S3 20,7 bp; S4 69,1 bp (59--79) & \etiketT{} Faz 2A ile tutarlı (73 bp) \\
Arketip dağılımı (kural + 10--15 düzeltmesi) & A 55, --- 49, B 39, G 31, C 30, D 17, F 8, E 1 & \etiketH{} \kod{veri/\allowbreak{}nace\_\allowbreak{}matris.\allowbreak{}csv}; diğer parçalarda düzeltilecek \\
Düzeltmeler & D-1 KN8 → NACE ağırlıkları; D-2 KKTC akımları; D-3 ``\%5'i aşan 21 sınıf'' yerine 9 (P50) / 17 (P10) & \kod{nace\_\allowbreak{}kartlar\_\allowbreak{}10-\allowbreak{}12.\allowbreak{}md} Ek D \\
\end{longtable}

\endgroup

\hypertarget{defter-faz-3-tabanux131-paruxe7a-2-1315}{%
\section{Faz 3 tabanı (parça 2: 13--15)}\label{defter-faz-3-tabanux131-paruxe7a-2-1315}}

Ayrıntı: \belgebag{nace_kartlar_13-15}{}{\kod{nace\_\allowbreak{}kartlar\_\allowbreak{}13-\allowbreak{}15.\allowbreak{}md}}.

\begingroup\gbtablo\small

\begin{longtable}[]{@{}
  >{\raggedright\arraybackslash}p{(\columnwidth - 4\tabcolsep) * \real{0.1726}}
  >{\raggedright\arraybackslash}p{(\columnwidth - 4\tabcolsep) * \real{0.6521}}
  >{\raggedright\arraybackslash}p{(\columnwidth - 4\tabcolsep) * \real{0.1753}}@{}}
\toprule\noalign{}
\begin{minipage}[b]{\linewidth}\raggedright
\textbf{Gösterge}
\end{minipage} & \begin{minipage}[b]{\linewidth}\raggedright
\textbf{Değer}
\end{minipage} & \begin{minipage}[b]{\linewidth}\raggedright
\textbf{Kaynak ve etiket}
\end{minipage} \\
\midrule\noalign{}
\endhead
\bottomrule\noalign{}
\endlastfoot
13--15 ağırlığı & İstihdam 1.238 bin (imalatın \%24,6'sı); KD \$22,5 mlr; ihracat \$35,9 mlr (\%14,3); AB'ye ihracat \$17,3 mlr (\%16,4) & \etiketV{} YSHİ 2024, BACI 2024 \\
AB'ye ihracatın rejimi ve S4 tarifesi & GB \%100; S4'te AB tarifesi \%10,2 (tekstil \%7,3, giyim \%11,8, deri \%7,4); giyimde AB'ye ihracatın \%98'i ≥ \%10 tarifeli satırlarda & \etiketV{}/\etiketT{} Comext 2025, KN 2026 \\
13--15, grup toplamı P50 (P10 / P90), T+5 & S2: AB'ye ihracat −\$1.854 mn (−4.010 / +347), istihdam −14.419 (−30.880 / +1.234), cari denge −\$893 mn. S4: AB'ye ihracat −\$4.104 mn (−6.394 / −1.698), istihdam −39.397 (−60.851 / −20.019), cari denge −\$2.634 mn (−3.841 / −1.473) & \etiketT{} \kod{veri/\allowbreak{}f3\_\allowbreak{}senaryo\_\allowbreak{}grup.\allowbreak{}csv} \\
PEM kümülasyonu (S4) & Ortaklara kumaş ve teknik tekstil ihracatı \$1.736 mn; maruziyet \$1.131 mn (Fas 613); S4 kaybı −\$673 mn (−858 / −483), istihdam −4.750 & \etiketT{} \kod{veri/\allowbreak{}f3\_\allowbreak{}pem\_\allowbreak{}ortak\_\allowbreak{}ab\_\allowbreak{}payi.\allowbreak{}csv}; V-40 \etiketH{} \\
TÜFE (birinci tur, P50) & S2 1,7 bp; S4 15,6 bp (dış giyim 4,3, ayakkabı 3,9); KDV takasıyla S4 3,6 bp & \etiketT{} f3e; Faz 2A Tablo 1.5a \\
Bölgeler (S4 P50) & İstanbul −12.209; Mardin--Batman--Şırnak--Siirt −3.388 (13--15 bölge imalatının \%78'i); Bursa −2.873; Tekirdağ −2.341; Gaziantep −2.228 & \etiketT{} f3b bölge payları \\
Duruşlar & Geliştir 11, Dönüştür 2 (14.13, 15.11), Yönetilen uyum 3 (14.14, 14.19, 14.39), Koru 5 & \etiketH{} \\
Düzeltmeler & D-4 EUDR (Tüzük 2026/2102: sığır derisi kapsam dışı); D-5 10--12'nin S2 sayıları; D-6 Faz 2'nin giyim tahmini (−59 bin) geçersiz & \kod{nace\_\allowbreak{}kartlar\_\allowbreak{}13-\allowbreak{}15.\allowbreak{}md} Ek D \\
\end{longtable}

\endgroup
